\section{Discussion: What the Rule Costs}
\label{sec:discussion}

beni's design follows a principle that a restriction belongs in the language only if it buys a
guarantee. This rule buys a \emph{performance} guarantee---every list write is $O(1)$, or the $O(n)$ a
hand-written program also pays; there is no hidden copy, no trie, no reference count, and a smaller
runtime---at the price of rejecting programs that are \emph{correct} under value semantics. This section
states that price without softening it.

\subsection{Programs that are rejected}

Every one of the following is a correct beni program today.
\begin{enumerate}
\item Keeping an old version: an undo history, a snapshot for a diff, a ``previous'' field in the model,
a test that holds its input and compares it with the output (\S\ref{sec:hard-undo}).
\item A command, fiber, subscription or stored closure that captured a list which a later message edits
(\S\ref{sec:hard-fibers})---including the beginner's \code{Cmd.perform (λsend → send (Saved model.todos))}.
\item A closure that edits a captured list and is passed to \code{map}, \code{filter}, \code{sortBy} or any
other $\mathit{many}$ position (\S\ref{sec:hard-closures}).
\item A list read out of a container that is still live: a dictionary lookup followed by a push, or a
push on an element of a list of lists while the outer list lives (\S\ref{sec:hard-containers}).
\item A list from JavaScript, a \code{Ref}, a \code{Queue}, or a \code{foreign} without a summary
(\S\ref{sec:hard-foreign}, \S\ref{sec:hard-ref}): a \emph{limit}, not sharing, and the message says so.
\item Anything the checker's joins or path cut cannot follow (\S\ref{sec:hard-branches}): limits.
\end{enumerate}
The fix is one call, \code{List.copy}, in every case; as Henriksen puts it for Futhark, ``an
alias-related type error can always be fixed by copying'' \citep{henriksen2022uniqueness}. What the fix
costs is a copy where the program meant to share. For the first case that is the program's intended
semantics made explicit; for the last two it is a copy the programmer pays for the checker's blindness.

\paragraph{The representation.} The trie's reason to exist is gone, and so is its persistence: a
program that wants versions pays $O(n)$ per version, explicitly. Prepending onto a long unique list is
\code{unshift}, $O(n)$, so an Elm-style prepend accumulator is quadratic again, unless the backend's
rewrite of building loops catches it; that rewrite catches a tail-recursive accumulator but not a
\code{foldr} lambda.

\paragraph{The API.} A public function's summary is part of its contract. Editing a body can break a
caller in another module with an ownership error that the caller's author did not write. The interface
firewall makes this visible, not painless.

\paragraph{Learners.} The error will fire first in \code{update}. Crichton and Krishnamurthi found the
chapter on ownership to be ``a significant drop-off point'' for readers of the Rust book
\citep[\S3.1]{crichton2024profiling}, and learners who can say why a program was rejected fix it less
than half the time \citep{crichton2023grounded}. The message design of \S\ref{sec:errors} is the
mitigation; it is not a guarantee.

\paragraph{What two real applications say.} We examined two applications written in beni: a TodoMVC
implementation (in four builds) and an implementation of Conduit, the RealWorld example application. By
our count neither calls \code{set}, \code{push}, \code{update}, \code{swap}, \code{pop}, \code{insertAt} or
\code{removeAt}; they edit lists with \code{List.map}, \code{List.filter} and \code{++}. Under the
nine-writer demand set, their only demands are their list concatenations. TodoMVC has one per build,
\code{model.todos ++ [ new ]}, where the record built just before it holds the old array at \code{.todos}
until a helper replaces that field: accepted. Conduit has sixteen. Seven have a model field on the left
(\code{\{ model | errors = model.errors ++ Api.errorMessages error \}}) and are accepted by the same
argument; the other nine have a list literal, or a list built earlier in the same function, on the left.
The core library has seven calls of \code{List.push}, every one an accumulator in a fold or a loop.

So the rule as stated changes nothing in these two applications. The in-place writes they would
benefit from are their \code{List.map model.todos (λt → …)} and \code{List.filter} calls, which are not
demands unless the demand set is widened (\S\ref{sec:disc-options}). If it is, TodoMVC's toggle becomes
\code{a[i] = f(a[i])} over a unique \code{todos}, and its delete compacts in place. But the same widened
rule also rejects \code{List.map} over a list the program keeps, and TodoMVC's view does keep one: its
filtered list is a derived value recomputed while the page holds the model (\S\ref{sec:hard-tea}).

\paragraph{Escape hatches.} \code{List.copy}, always. A persistent-vector library written in beni for
programs that want versions. And a \emph{warning mode} that keeps the persistent representation and
copies where the checker would reject, printing the same message as a warning. That is the form every
language with a run-time fallback ships in some way---Lean offers \code{dbgTraceIfShared}, which prints
a message ``if the argument's reference count is greater than one'' \citep{lean2025rc}---and it gives
up the one thing the rule is for: an accepted program's cost would again depend on whether its author
read a warning. We do not propose it as the end state; we propose measuring first
(\S\ref{sec:evaluation}) and deciding between error and warning on the numbers.

\subsection{Design options}
\label{sec:disc-options}

The language specification can change where the checker needs it. Each option below says what it buys
and what it costs; none is required for soundness.

\paragraph{Widen the demand set to same-shape rebuilds.} Make \code{map}, \code{indexedMap},
\code{filter}, \code{filterMap}, \code{reverse}, \code{sort}, \code{sortBy}, \code{sortWith}, \code{take}
and \code{slice} consuming writers: \code{map} becomes \code{a[i] = f(a[i])}, \code{filter} compacts in
place, \code{sort} is \code{a.sort(cmp)} (stable in every current engine), \code{take n} is
\code{a.length = n}. This is where real programs edit lists, and it would make nested edits through
\code{map} acceptable when contents are exclusive. It also rejects every read-after-\code{map} in
\code{update}, and it meets a wall: the same \code{List.filter model.todos …} appears in \code{view},
where it is a derived value that must build fresh. The option is therefore either two functions (an
in-place and a fresh \code{map}---the code duplication FP\textsuperscript{2} identifies as the cost of
static uniqueness typing, ``where a single function can have multiple different implementations''
\citep[\S1.4.1]{lorenzen2023fp2}) or one spelling whose meaning is decided by position: a demand in an
\code{update} arm, a fresh build in \code{view} or a derived value. We would adopt it only with the
evaluation's numbers in hand, since this is where most false errors would come from, and only in the
by-position form.

\paragraph{Drop the identity promises that alias.} beni's library currently promises that
\code{filter}, \code{map}, \code{slice}, \code{take}, \code{xs ++ []}, \code{[] ++ ys} and a \code{concat}
with one non-empty list return an input itself in some cases. Under the rule a result that may be the
input names it in $\mathit{res}$, so \code{xs ▷ List.filter p ▷ List.push y} demands \code{xs} unique
although \code{filter} usually returns a fresh array; and \code{[] ++ ys} returning \code{ys} would hand
out \code{ys}'s cell from a consumed \code{append}. Dropping these promises makes those results always
fresh (or in place, under the previous option). The renderer loses one pointer test for a
\code{filter} that kept everything. The promises for out-of-range \code{set}, \code{swap} and
\code{removeAt} stay: they are no-ops on the unique array. (The library's specification is itself
inconsistent about \code{filter}'s identity today, and must be settled before its $\mathit{res}$ can be
declared.)

\paragraph{\code{List.copy}.} New, shallow, fresh: the explicit escape hatch and the mechanical fix. A
copy of a list the checker can see is unique could be a warning, never an error.

\paragraph{\code{++} on lists consumes its left operand.} \code{xs ++ ys} pushes each element of
\code{ys} onto \code{xs}, consistently with \code{[ …xs, …ys ]}; \code{xs ++ ys} with \code{xs} read later
becomes an error. A smaller version of the first option.

\paragraph{Ownership words on \code{foreign} declarations, optionally on public annotations.} A
\code{foreign} parameter may say consumed, borrowed or shared; a result fresh or ``is parameter $i$''.
This is required for soundness at the boundary (A2, A6). Core and the platforms may assert a summary
their bodies cannot show; every other package may only \emph{narrow} an inferred one. A public beni
declaration may optionally state its summary, which must then be at most what the body takes: a stated
borrowed where the body consumes is an error at the declaration, while a stated consumed where the body
borrows is accepted and freezes the interface at consumed. This answers interface churn for library
authors who want stable signatures---at Jane Street, OxCaml interface files carried 27\,382 locality
annotations, because ``the implementations can generally infer locality, whereas interfaces must write
it down explicitly''
\citep[\S7]{lorenzen2024oxidizing}---without asking anything of application code. The vocabulary must
also reach a function type inside a platform's record, where the hand-over obligation is written.

\paragraph{Closures that consume a capture are \code{once}.} Already a consequence of the rules, listed
because it is visible in the language: \code{List.map xs (λx → List.push acc x)} is an error whose fix is
a fold. It stays an error, because two calls would write one array.

\paragraph{A persistent vector as a library.} The current trie, written in beni with no in-place
writes, as a library module for programs that keep versions---the status of a dictionary, not of
\code{List}.

\paragraph{Headroom for prepend.} \code{[ x, …xs ]} on a unique plain array is \code{unshift}. A view
with headroom ($o > 0$) could prepend in $O(1)$ by writing $b[o - 1]$, at the cost of a growth policy in
the list runtime. To be decided on the evaluation's count of prepends outside building loops.

\paragraph{Sending consumes its message.} A fiber that sends a list and reads it afterwards is an error
at the sender, so payloads are unique at the arm; the alternative is every payload shared at every arm.

\paragraph{Writing the dependence down.} The rule depends on the language's evaluation order being the
order the emitted code runs in. It already is, and A5 keeps it so; the language definition should say
that in-place writes are calls for the purpose of ordering.

\subsection{Limitations and open questions}
\label{sec:disc-open}

\begin{enumerate}
\item \emph{The rejection rate on real code}, exact and approximate, for both demand sets. Nothing in the
literature we know of stands in for it; \S\ref{sec:evaluation} proposes how to measure it.
\item \emph{Mutable cells} (\S\ref{sec:hard-ref}): a sound discipline for in-place edits of a cell's
contents needs a whole-program argument about readers; it is left as a limit with a copy as the fix.
\item \emph{Path-sensitive cell sets} (\S\ref{sec:hard-branches}): worth their cost only if the join tag is
frequent.
\item \emph{Path-insensitive use}: \code{f xs flag = if flag then List.length xs else List.length
(List.push xs 1)} is summarised as consuming \code{xs}, so \code{n = f xs True; List.length xs} is a false
error. A summary conditional on a \emph{value}, rather than on a class, is outside the summary grammar.
\item \emph{Interface churn} under inferred summaries, and the warm rebuild budget, since a body edit
that changes a summary is a signature edit.
\item \emph{Exclusive contents in practice}: how often the containers that real programs edit through
their elements are built exclusively, and which construction most often clears the flag. If decoded
data is not declared element-fresh, every nested edit of loaded data is a limit.
\item \emph{The two-slots-one-record hazard} is the one place the design relies on a flag that the
program's construction must establish rather than on liveness alone. We would direct a reviewer to try
to break Lemma~\ref{lem:excl} with containers built through constrained generic code and through
\code{foreign} siblings.
\item \emph{The entry fixpoint's cost}: one whole-program pass over \code{init} and the arms per build,
the only non-modular computation in the design.
\item \emph{A \textit{once} closure called once in fact} but passed to a $\mathit{many}$ position stays
an error.
\item \emph{The development crash screen} must print no program value (A9).
\item \emph{Primary sources we could not check}: the best-typing result that Aspinall et al.\ attribute
to Kone\v{c}n\'y, and a formal statement of the ``two semantics coincide'' theorem for Cogent, should be
checked against the originals before the design is specified.
\end{enumerate}
